\chapter{System Design}

\section{System Architecture}
The system has a layered architecture (Fig.~\ref{fig:arch}). The web interface passes an upload to the validation layer, which reads the true file type from the bytes, checks the size and decodes the file. A router then calls the image, video or document analyser. Every analyser implements the same interface and returns an \emph{evidence record}: a score, a visual (heatmap, timeline or mask), numeric cues and flags raised by rules. The scoring layer calibrates the record into a probability and a Trust Score, and the explanation layer writes the sentences. Because the analysers only share this interface, a model can be replaced without touching the application, and the three phases of the project can be built one after the other.

\begin{figure}[htbp]
  \centering
  % Layered architecture: solid = implemented, dashed = designed
\resizebox{\linewidth}{!}{%
\begin{tikzpicture}
\node[box,text width=3.8cm] (up) at (1.9,0) {User upload\\image, video or document};
\node[box,text width=5.4cm] (val) at (7.6,0) {Validation and routing\\file type read from the bytes,\\size limit, safe decoding};
\draw[arr] (up) -- (val);
% bus to the three analysers
\draw[thick,cLine] (val.south) -- (7.6,-1.25);
\draw[thick,cLine] (1.9,-1.25) -- (13.3,-1.25);
\node[boxg,text width=4.1cm] (img) at (1.9,-3.0) {\textbf{Image analyser}\\(implemented)\\native-resolution patches,\\frozen CLIP + trained head,\\patch heatmap};
\node[designed,text width=4.1cm] (vid) at (7.6,-3.0) {\textbf{Video analyser}\\(designed)\\sampled frames, face crops,\\frame scores, timeline};
\node[designed,text width=4.1cm] (doc) at (13.3,-3.0) {\textbf{Document analyser}\\(designed)\\OCR, ELA / SRM / DCT views,\\dual-stream U-Net mask,\\rule checks};
\draw[arr] (1.9,-1.25) -- (img.north);
\draw[arr] (7.6,-1.25) -- (vid.north);
\draw[arr] (13.3,-1.25) -- (doc.north);
% evidence bus
\draw[thick,cLine] (img.south) -- (1.9,-5.0);
\draw[thick,cLine] (vid.south) -- (7.6,-5.0);
\draw[thick,cLine] (doc.south) -- (13.3,-5.0);
\draw[thick,cLine] (1.9,-5.0) -- (13.3,-5.0);
\node[box,text width=6.4cm] (ev) at (7.6,-6.1) {Evidence records\\score, map or timeline, numeric cues, flags};
\draw[arr] (7.6,-5.0) -- (ev.north);
\node[box,text width=6.4cm] (cal) at (7.6,-7.7) {Calibration and fusion\\$p=\sigma(a\,s+c)$ (Platt scaling)};
\draw[arr] (ev) -- (cal);
\node[boxo,text width=8.0cm] (out) at (7.6,-9.3) {Trust Score 0--100, verdict band,\\explanation sentences, downloadable report};
\draw[arr] (cal) -- (out);
\end{tikzpicture}}

  \caption{Layered architecture. Solid boxes are implemented; dashed boxes are designed.}\label{fig:arch}
\end{figure}

\section{Training and Inference Are Separate}
Heavy computation, such as extracting features from hundreds of thousands of patches, runs once on a free cloud GPU. The application loads the frozen encoder and a trained head of a few kilobytes and runs on a CPU. Loading only the image half of CLIP, tensor by tensor, keeps the peak memory near 1.5~GB. Fig.~\ref{fig:train} shows both sides.

\begin{figure}[htbp]
  \centering
  % Training (offline) and inference (online) workflow
\resizebox{\linewidth}{!}{%
\begin{tikzpicture}
\node[font=\small\bfseries] at (7.2,1.0) {Offline: free cloud GPU and laptop};
\node[box,text width=2.9cm] (a1) at (0,0) {Public datasets\\(OpenFake, DiffSeg30k)};
\node[box,text width=2.9cm] (a2) at (3.7,0) {Patches at native size + degraded copies};
\node[box,text width=2.9cm] (a3) at (7.4,0) {Frozen CLIP features, stored once};
\node[box,text width=2.9cm] (a4) at (11.1,0) {Fit small head, choose $C$ on validation};
\node[box,text width=2.9cm] (a5) at (14.8,0) {Platt scaling and verdict limits on validation};
\draw[arr] (a1) -- (a2);
\draw[arr] (a2) -- (a3);
\draw[arr] (a3) -- (a4);
\draw[arr] (a4) -- (a5);
\node[boxw,text width=3.6cm] (t) at (14.8,-2.2) {Held-out test split scored once};
\draw[arr] (a5) -- (t);
\node[store,text width=3.6cm] (w) at (9.2,-2.2) {Weight files\\(a few kilobytes)};
\draw[arr] (t) -- (w);
\node[font=\small\bfseries] at (7.2,-3.6) {Online: CPU only};
\node[boxg,text width=2.9cm] (b1) at (0,-5.0) {Upload and validate};
\node[boxg,text width=2.9cm] (b2) at (3.7,-5.0) {Ten patches, JPEG q95};
\node[boxg,text width=2.9cm] (b3) at (7.4,-5.0) {CLIP encoder + loaded head};
\node[boxg,text width=2.9cm] (b4) at (11.1,-5.0) {Calibrated Trust Score};
\node[boxo,text width=2.9cm] (b5) at (14.8,-5.0) {Verdict, heatmap, explanation};
\draw[arr] (b1) -- (b2);
\draw[arr] (b2) -- (b3);
\draw[arr] (b3) -- (b4);
\draw[arr] (b4) -- (b5);
\draw[darr] (w.south) -- ++(0,-1.3) -| (b3.north);
\end{tikzpicture}}

  \caption{Offline training workflow and online use of the trained model.}\label{fig:train}
\end{figure}

\section{Data Design}
Features of the frozen encoder are computed once and stored with their label, source (the generator), view type and the kind of patch. Because features are stored, a new head or a new calibration costs minutes instead of hours.

\begin{table}[htbp]
\centering\small
\caption{Data used for the image analyser.}
\begin{tabularx}{\linewidth}{@{}p{3.3cm}X@{}}
\toprule
Source & Use\\
\midrule
OpenFake \cite{openfake} & Real photographs and images from about sixty generators. Training, validation and test splits are made by generator; the test split contains generators that never appear in training (for example GPT-Image-1.5 and Nano Banana Pro)\\
DiffSeg30k \cite{diffseg} and similar edit data & Patches of locally edited photographs with masks, used as extra training patches (a share is held out for testing)\\
Degraded copies & Every training image also contributes a randomly resized, blurred and JPEG-compressed copy\\
Team photographs & Full-size phone photographs taken by the team, their edits made with ChatGPT and Gemini; used only for testing\\
\bottomrule
\end{tabularx}
\end{table}

For the video and document phases the sources are FaceForensics++ and Celeb-DF (video) and clean public documents with synthetic forgeries, DocTamper and human-made forgeries (documents). Splits follow the unit of the problem: by generator for images, by video for video, by document for documents.

\section{Design of the Image Analyser}
The image is not resized, because resizing removes fine traces. Ten patches of $224\times224$ pixels (the centre and a $3\times3$ grid spread over the image) are cut at native resolution. Each patch is re-saved as a JPEG with quality 95, so that real and generated inputs share one compression history and the file format cannot become a shortcut \cite{grommelt}. The frozen CLIP ViT-L/14 encoder $f$ maps a patch $x_i$ to a normalised feature $z_i=f(x_i)/\lVert f(x_i)\rVert$, and a linear head scores it: $s_i=\mathbf{w}^{\mathsf T}z_i+b$. Fig.~\ref{fig:pipe} shows the steps.

\begin{figure}[htbp]
  \centering
  % Image analyser pipeline (two rows)
\resizebox{\linewidth}{!}{%
\begin{tikzpicture}
\node[box,text width=3.0cm] (b1) at (0,0) {\textbf{1.} Input image\\(not resized)};
\node[box,text width=3.0cm] (b2) at (4.0,0) {\textbf{2.} Ten patches\\224$\times$224: centre\\+ 3$\times$3 grid};
\node[box,text width=3.0cm] (b3) at (8.0,0) {\textbf{3.} Re-save each patch\\as JPEG, quality 95};
\node[box,text width=3.0cm] (b4) at (12.0,0) {\textbf{4.} Frozen CLIP\\ViT-L/14 encoder,\\$z=f(x)/\lVert f(x)\rVert$};
\draw[arr] (b1) -- (b2);
\draw[arr] (b2) -- (b3);
\draw[arr] (b3) -- (b4);
\node[boxg,text width=3.0cm] (b5) at (0,-3.0) {\textbf{5.} Linear head\\$s=\mathbf{w}^{\mathsf T}z+b$\\per patch};
\node[boxg,text width=3.0cm] (b6) at (4.0,-3.0) {\textbf{6.} Image score\\= mean of the ten\\patch scores};
\node[boxg,text width=3.0cm] (b7) at (8.0,-3.0) {\textbf{7.} Platt scaling\\$p=\sigma(a\,s+c)$};
\node[boxo,text width=3.0cm] (b8) at (12.0,-3.0) {\textbf{8.} Trust Score\\$T=\mathrm{round}(100(1-p))$\\and verdict band};
\draw[arr] (b4.south) -- ++(0,-0.75) -| (b5.north);
\draw[arr] (b5) -- (b6);
\draw[arr] (b6) -- (b7);
\draw[arr] (b7) -- (b8);
\node[font=\footnotesize,text width=3.0cm,align=center] at (4.0,-4.6) {patch scores also form the $3\times3$ heatmap};
\draw[darr] (4.0,-4.25) -- (4.0,-3.55);
\end{tikzpicture}}

  \caption{Processing steps of the image analyser.}\label{fig:pipe}
\end{figure}

\subsection{Mathematical Model}
Let $z_1,\dots,z_{10}$ be the normalised features of the ten patches of one image.
\begin{align}
  s_i &= \mathbf{w}^{\mathsf T} z_i + b, & s &= \frac{1}{10}\sum_{i=1}^{10} s_i, \\
  p &= \sigma(a\,s + c) = \frac{1}{1+e^{-(a s + c)}}, & T &= \mathrm{round}\bigl(100\,(1-p)\bigr).
\end{align}
Only $\mathbf{w}$ and $b$ are trained, by L2-regularised logistic regression with class weights; the penalty strength $C$ is chosen on a validation split. The Platt parameters $a$ and $c$ are fitted on validation data, so that $p$ behaves like a probability. The verdict limits are chosen on validation data, too: at most 3\% of real validation images may be called AI-generated and at most 10\% of AI images may be called authentic. For the image analyser they are $T\ge 54$ (Likely Authentic) and $T\le 17$ (Likely AI-generated / Manipulated); the values in between mean \emph{Needs Human Review}.

\section{Designs of the Video and Document Analysers}
\textbf{Video analyser (designed).} About one frame per second is sampled (at most 32 frames). MTCNN \cite{mtcnn} finds and aligns faces, and every crop keeps a margin so that the boundary of the face stays visible \cite{li}. Each crop is scored, first by the image analyser unchanged and then by a head trained on FaceForensics++ \cite{rossler}. The frame scores are averaged over the top-scoring share of frames, so that a short manipulated segment is not diluted by clean frames; they also give a timeline. Splits are made by video, never by frame, and transfer is tested on Celeb-DF \cite{celebdf}.

\textbf{Document analyser (designed).} Pages are rendered and read by OCR \cite{tesseract}. Error-level analysis \cite{krawetz}, spatial-rich-model residuals \cite{fridrich} and DCT features are stacked as a second input stream next to the RGB page of a dual-stream U-Net \cite{unet,zhou} with an EfficientNet encoder \cite{effnet}. It is trained with binary cross-entropy plus a Dice term \cite{vnet} on synthetic forgeries (text replacement, copy-move, splicing) with exact masks and scanner-like noise, and tested on human-made forgeries and DocTamper \cite{doctamper}. Rule checks (line items against the total, dates, font heights, PDF modification history) add flags that are auxiliary evidence and never proof.

\section{Data Flow Diagrams}
The data flow diagrams show how data moves through the system at three levels of detail. External entities are drawn as filled rectangles, processes as circles, data stores as grey boxes, and every arrow is labelled with the data it carries.

\subsection{DFD Level 0 (Context Diagram)}
At level 0 the whole system is one process (Fig.~\ref{fig:dfd0}). The user sends an image file and the chosen detector model and receives the Trust Score, verdict, heatmap, explanation, notices and report, or an error message. The project team supplies training data and settings and receives evaluation results and calibration reports. Public datasets supply labelled real and AI-generated images, and a model hub supplies the pretrained CLIP encoder weights.

\begin{figure}[htbp]
  \centering
  % Data flow diagram, level 0 (context diagram): the whole system as one process and every external entity
\resizebox{\linewidth}{!}{%
\begin{tikzpicture}
\node[proc,minimum size=3.6cm,text width=2.9cm,font=\small] (sys) at (0,0) {\textbf{0}\\Deepfake Detection System};
\node[actor,text width=2.6cm] (user) at (-7.5,0) {\textbf{User}\\(student, journalist, analyst)};
\node[actor,text width=2.6cm] (team) at (7.5,0) {\textbf{Project team}\\(administrator)};
\node[actor,text width=3.0cm] (data) at (-3.6,-4.4) {\textbf{Public datasets}\\(OpenFake, DiffSeg30k)};
\node[actor,text width=3.0cm] (clip) at (3.6,-4.4) {\textbf{Model hub}\\(pretrained CLIP weights)};
\draw[arr] ($(user.east)+(0,0.45)$) -- node[above,font=\scriptsize] {image file, model choice} ($(sys.west)+(0,0.45)$);
\draw[arr] ($(sys.west)+(0,-0.45)$) -- node[below,font=\scriptsize,align=center] {Trust Score, verdict, heatmap,\\explanation, notices, report,\\or error message} ($(user.east)+(0,-0.45)$);
\draw[arr] ($(team.west)+(0,0.45)$) -- node[above,font=\scriptsize] {training data, settings} ($(sys.east)+(0,0.45)$);
\draw[arr] ($(sys.east)+(0,-0.45)$) -- node[below,font=\scriptsize,align=center] {evaluation results,\\calibration report} ($(team.west)+(0,-0.45)$);
\draw[arr] (data.north) -- node[left,font=\scriptsize,align=center,pos=0.45] {labelled real\\and AI images} (sys.south west);
\draw[arr] (clip.north) -- node[right,font=\scriptsize,align=center,pos=0.45] {encoder\\weights} (sys.south east);
\end{tikzpicture}}

  \caption{Data flow diagram, level 0 (context diagram).}\label{fig:dfd0}
\end{figure}

\subsection{DFD Level 1}
Level 1 opens the system into six processes and five data stores (Fig.~\ref{fig:dfd1}).
\begin{itemize}
  \item \textbf{1.0 Validate and decode upload} reads the file type from the bytes, checks the size and decodes the image; an invalid file goes back to the user as an error message.
  \item \textbf{2.0 Analyse image} cuts the patches, encodes them with the CLIP encoder (D1) and scores them with the trained head (D2). It sends the image score to 3.0 and the patch scores to 5.0 for the heatmap.
  \item \textbf{3.0 Calibrate and assign verdict} turns the score into a probability with the calibration (D3) and assigns the band.
  \item \textbf{4.0 Check file type and provenance} uses the watermark template and notice rules (D4) to detect drawings, small or recompressed files and the visible Gemini watermark, and reads the camera and C2PA metadata.
  \item \textbf{5.0 Explain and build report} combines the verdict, the patch scores and the notices into the result card, the heatmap, the explanation and the report for the user.
  \item \textbf{6.0 Train and calibrate model} (offline, run by the project team) takes labelled images from the public datasets, stores patch features (D5), trains a new head (D2), chooses a new calibration (D3) and returns the evaluation results.
\end{itemize}

\begin{figure}[htbp]
  \centering
  % Data flow diagram, level 1
\resizebox{\linewidth}{!}{%
\begin{tikzpicture}
\node[actor,text width=2.0cm] (user) at (0,0) {User};
\node[proc,minimum size=2.6cm,text width=2.1cm] (p1) at (4.0,0) {\textbf{1.0}\\Validate upload};
\node[proc,minimum size=2.6cm,text width=2.1cm] (p2) at (8.2,0) {\textbf{2.0}\\Analyse with the routed analyser};
\node[proc,minimum size=2.6cm,text width=2.1cm] (p3) at (12.4,0) {\textbf{3.0}\\Calibrate and score};
\node[proc,minimum size=2.6cm,text width=2.1cm] (p4) at (12.4,-3.8) {\textbf{4.0}\\Explain and report};
\node[store,text width=3.1cm] (d1) at (8.2,-2.6) {D1  Trained heads\\and CLIP encoder};
\node[store,text width=3.1cm] (d2) at (8.2,-4.6) {D2  Calibration and\\verdict limits};
\draw[arr] (user) -- node[above,font=\scriptsize] {file} (p1);
\draw[arr] (p1) -- node[above,font=\scriptsize] {valid file} (p2);
\draw[arr] (p2) -- node[above,font=\scriptsize] {evidence} (p3);
\draw[arr] (p3) -- node[right,font=\scriptsize] {$T$, verdict} (p4);
\draw[arr] (d1.north) -- (p2.south);
\draw[arr] (d2.east) -- (p3.south west);
\draw[arr] (p1.north) -- ++(0,1.5) node[above,font=\scriptsize] {error message} -| (user.north);
\draw[arr] (p4.south) -- ++(0,-1.3) -| (user.south);
\node[font=\scriptsize,align=center] at (6.2,-6.6) {Trust Score, verdict, heatmap, explanation, report};
\end{tikzpicture}}

  \caption{Data flow diagram, level 1.}\label{fig:dfd1}
\end{figure}

\subsection{DFD Level 2}
Level 2 opens the two central processes (Fig.~\ref{fig:dfd2}). Process 2.0 consists of 2.1 cutting ten patches at native resolution, 2.2 re-saving each patch as JPEG quality 95, 2.3 encoding the patches with CLIP (D1), 2.4 scoring each patch with the head (D2) and 2.5 averaging the scores and building the $3\times3$ grid for the heatmap. Process 3.0 consists of 3.1 Platt scaling with the parameters $a$ and $c$ (D3), 3.2 computing the Trust Score $T=\mathrm{round}(100(1-p))$ and 3.3 assigning the verdict band with the limits 54 and 17 (D3).

\begin{figure}[htbp]
  \centering
  % Data flow diagram, level 2: processes 2.0 (Analyse image) and 3.0 (Calibrate and assign verdict) opened up
\resizebox{\linewidth}{!}{%
\begin{tikzpicture}[sp/.style={proc,minimum size=2.5cm,text width=2.05cm,font=\footnotesize},ds/.style={store,text width=2.4cm,font=\footnotesize}]
% ---- process 2.0
\node[sp] (a1) at (0,0) {\textbf{2.1}\\Cut 10 patches, native resolution};
\node[sp] (a2) at (4.2,0) {\textbf{2.2}\\Re-save each patch as JPEG q95};
\node[sp] (a3) at (8.4,0) {\textbf{2.3}\\Encode patches with CLIP};
\node[sp] (a4) at (12.6,0) {\textbf{2.4}\\Score each patch with the head};
\node[sp] (a5) at (16.8,0) {\textbf{2.5}\\Average scores, build $3\times3$ grid};
\node[ds] (d1) at (8.4,3.2) {D1 \;CLIP encoder weights};
\node[ds] (d2) at (12.6,3.2) {D2 \;Trained heads};
% ---- process 3.0
\node[sp] (b1) at (8.4,-5.6) {\textbf{3.1}\\Platt scaling $p=\sigma(a\,s+c)$};
\node[sp] (b2) at (12.6,-5.6) {\textbf{3.2}\\Compute Trust Score $T$};
\node[sp] (b3) at (16.8,-5.6) {\textbf{3.3}\\Assign verdict band};
\node[ds] (d3) at (12.6,-9.0) {D3 \;Calibration and verdict limits};
% ---- boundaries of the two level-1 processes
\begin{scope}[on background layer]
\node[draw=cLine,dashed,thick,rounded corners=6pt,fit=(a1)(a5)(d1)(d2),inner sep=10pt,label={[font=\small\bfseries,anchor=north west]north west:Process 2.0: Analyse image}] {};
\node[draw=cLine,dashed,thick,rounded corners=6pt,fit=(b1)(b3)(d3),inner sep=10pt,label={[font=\small\bfseries,anchor=north east]north east:Process 3.0: Calibrate and assign verdict}] {};
\end{scope}
% ---- in and out
\node[boxw,text width=2.6cm,font=\footnotesize] (in) at (-4.4,0) {from 1.0\\decoded image};
\node[boxw,text width=2.6cm,font=\footnotesize] (outh) at (21.4,0) {to 5.0\\patch scores for heatmap};
\node[boxw,text width=2.6cm,font=\footnotesize] (outv) at (21.4,-5.6) {to 5.0\\Trust Score, verdict};
% ---- flows
\draw[arr] (in) -- (a1);
\draw[arr] (a1) -- node[above,font=\scriptsize] {10 patches} (a2);
\draw[arr] (a2) -- node[above,font=\scriptsize] {JPEG patches} (a3);
\draw[arr] (a3) -- node[above,font=\scriptsize] {features $z_i$} (a4);
\draw[arr] (a4) -- node[above,font=\scriptsize] {scores $s_i$} (a5);
\draw[arr] (d1) -- node[right,font=\scriptsize] {encoder} (a3);
\draw[arr] (d2) -- node[right,font=\scriptsize] {$\mathbf{w}, b$} (a4);
\draw[arr] (a5) -- node[above,font=\scriptsize] {grid} (outh);
\draw[arr] (a5.south) -- ++(0,-2.0) -| node[pos=0.25,above,font=\scriptsize] {image score $s$} (b1.north);
\draw[arr] (b1) -- node[above,font=\scriptsize] {$p$} (b2);
\draw[arr] (b2) -- node[above,font=\scriptsize] {$T$} (b3);
\draw[arr] (b3) -- (outv);
\draw[arr] (d3.west) -| node[pos=0.3,below,font=\scriptsize] {$a, c$} (b1.south);
\draw[arr] (d3.east) -| node[pos=0.3,below,font=\scriptsize] {limits 54 and 17} (b3.south);
\end{tikzpicture}}

  \caption{Data flow diagram, level 2: processes 2.0 and 3.0.}\label{fig:dfd2}
\end{figure}

\section{Use-Case Diagram}
The use-case diagram (Fig.~\ref{fig:uc}) shows every function of the system and the three actors that use it.
\begin{itemize}
  \item \textbf{User:} uploads an image or tries a sample image, chooses the detector model, views the Trust Score and verdict, the patch heatmap, the explanation, the metadata, frequency spectrum and error-level analysis, views the cross-generator results and (planned) downloads a report.
  \item \textbf{Included use cases:} uploading always includes validating the file and analysing the image; viewing the Trust Score always includes calibrating the score and assigning the band.
  \item \textbf{Extending use cases:} a notice for drawings, small or compressed files and the Gemini watermark alert extend the verdict only when they apply; the finer $6\times6$ heatmap extends the heatmap when the user asks for it.
  \item \textbf{Project team (administrator):} collects labelled datasets, extracts patch features on a cloud GPU, trains the classifier head, calibrates and sets the verdict limits, evaluates on unseen generators, updates the watermark and notice rules and deploys the application.
  \item \textbf{Hosting platform (Streamlit Cloud):} takes part in deploying the application.
\end{itemize}

\begin{figure}[htbp]
  \centering
  % Use-case diagram (UML): every function of the system, with include / extend relations and three actors
\resizebox{\linewidth}{!}{%
\begin{tikzpicture}[
  uc/.style={usecase,text width=3.3cm,minimum height=9mm},
  uca/.style={usecase,text width=3.3cm,minimum height=9mm,fill=cOrange},
  rel/.style={->,dashed,thick,draw=cLine},
  lab/.style={font=\scriptsize,fill=white,inner sep=1pt},
  stick/.pic={\draw[thick,cLine] (0,0.62) circle (0.24); \draw[thick,cLine] (0,0.38)--(0,-0.3) (-0.42,0.15)--(0.42,0.15) (0,-0.3)--(-0.32,-0.85) (0,-0.3)--(0.32,-0.85);}]
% ---- system boundary
\draw[thick,cLine,rounded corners=4pt,fill=white] (2.3,1.5) rectangle (19.4,-12.6);
\node[font=\bfseries] at (10.85,1.0) {Deepfake Detection System};
% ---- actors
\pic at (0,-5.4) {stick}; \node[font=\small\bfseries] at (0,-6.6) {User};
\pic at (21.9,-3.4) {stick}; \node[font=\small\bfseries,align=center] at (21.9,-4.75) {Project team\\(administrator)};
\pic at (21.9,-10.4) {stick}; \node[font=\small\bfseries,align=center] at (21.9,-11.75) {Hosting platform\\(Streamlit Cloud)};
% ---- user use cases (left column)
\node[uc] (u1) at (5.0,-0.2) {Upload an image};
\node[uc] (u2) at (5.0,-1.6) {Try a sample image};
\node[uc] (u3) at (5.0,-3.0) {Choose detector model};
\node[uc] (u4) at (5.0,-4.4) {View Trust Score and verdict};
\node[uc] (u5) at (5.0,-5.8) {View patch heatmap};
\node[uc] (u6) at (5.0,-7.2) {Read the explanation};
\node[uc] (u7) at (5.0,-8.6) {View metadata, spectrum and ELA};
\node[uc] (u8) at (5.0,-10.0) {View cross-generator results};
\node[uc] (u9) at (5.0,-11.4) {Download report (planned)};
% ---- included / extending use cases (middle column, beside the use case they belong to)
\node[uc] (m1) at (10.9,-0.2) {Validate file type and size};
\node[uc] (m2) at (10.9,-1.6) {Analyse image (patches, CLIP, head)};
\node[uc] (m4) at (10.9,-3.0) {Show notice: drawing, small or compressed file};
\node[uc] (m3) at (10.9,-4.4) {Calibrate score, assign band};
\node[uc] (m5) at (10.9,-5.8) {Show Gemini watermark alert};
\node[uc] (m6) at (10.9,-7.2) {Show finer $6\times6$ heatmap};
% ---- administrator use cases (right column)
\node[uca] (a1) at (16.6,-0.2) {Collect labelled datasets};
\node[uca] (a2) at (16.6,-1.9) {Extract patch features (cloud GPU)};
\node[uca] (a3) at (16.6,-3.6) {Train classifier head};
\node[uca] (a4) at (16.6,-5.3) {Calibrate, set verdict limits};
\node[uca] (a5) at (16.6,-7.0) {Evaluate on unseen generators};
\node[uca] (a6) at (16.6,-8.7) {Update watermark and notice rules};
\node[uca] (a7) at (16.6,-10.4) {Deploy the application};
% ---- associations
\foreach \uu in {u1,u2,u3,u4,u5,u6,u7,u8,u9}{\draw[thick,cLine] (0.55,-5.1) -- (\uu.west);}
\foreach \aa in {a1,a2,a3,a4,a5,a6,a7}{\draw[thick,cLine] (21.35,-3.1) -- (\aa.east);}
\draw[thick,cLine] (21.35,-10.1) -- (a7.east);
% ---- include / extend (arrow from the base use case to the included one; from the extension to the base)
\draw[rel] (u1) -- node[lab] {\guillemotleft include\guillemotright} (m1);
\draw[rel] (u2) -- node[lab] {\guillemotleft include\guillemotright} (m2);
\draw[rel] (u1.south east) -- node[lab,pos=0.62] {\guillemotleft include\guillemotright} (m2.north west);
\draw[rel] (u4) -- node[lab] {\guillemotleft include\guillemotright} (m3);
\draw[rel] (m4.south west) -- node[lab,pos=0.42] {\guillemotleft extend\guillemotright} (u4.north east);
\draw[rel] (m5.north west) -- node[lab,pos=0.42] {\guillemotleft extend\guillemotright} (u4.south east);
\draw[rel] (m6.north west) -- node[lab,pos=0.42] {\guillemotleft extend\guillemotright} (u5.south east);
\draw[rel] (a3) -- node[lab,right] {\guillemotleft include\guillemotright} (a2);
\draw[rel] (a2) -- node[lab,right] {\guillemotleft include\guillemotright} (a1);
\draw[rel] (a5) -- node[lab,right] {\guillemotleft include\guillemotright} (a4);
\end{tikzpicture}}

  \caption{Use-case diagram.}\label{fig:uc}
\end{figure}

\section{Sequence Diagram}
Fig.~\ref{fig:seq} shows the messages exchanged during one image check, from the upload to the result card.

\begin{figure}[htbp]
  \centering
  % Sequence diagram of one image check
\resizebox{\linewidth}{!}{%
\begin{tikzpicture}
\node[actor,minimum width=2.4cm] (u) at (0,0) {User};
\node[box,minimum width=2.6cm,text width=2.4cm] (w) at (3.7,0) {Web interface (Streamlit)};
\node[box,minimum width=2.6cm,text width=2.4cm] (r) at (7.4,0) {Validation and router};
\node[boxg,minimum width=2.6cm,text width=2.4cm] (a) at (11.1,0) {Image analyser};
\node[boxo,minimum width=2.6cm,text width=2.4cm] (s) at (14.8,0) {Scoring and explanation};
\foreach \x in {0,3.7,7.4,11.1,14.8}{\draw[dashed,cLine] (\x,-0.7) -- (\x,-10.2);}
\draw[arr] (0,-1.3) -- node[above,font=\scriptsize] {upload file} (3.7,-1.3);
\draw[arr] (3.7,-2.5) -- node[above,font=\scriptsize] {file bytes} (7.4,-2.5);
\node[font=\scriptsize,fill=white,draw=cLine,rounded corners=2pt] at (7.4,-3.5) {read type, check size, decode};
\draw[arr] (7.4,-4.6) -- node[above,font=\scriptsize] {analyse(image)} (11.1,-4.6);
\node[font=\scriptsize,fill=white,draw=cLine,rounded corners=2pt,text width=3.2cm,align=center] at (11.1,-5.8) {10 patches, JPEG q95, CLIP features, linear head, mean};
\draw[darr] (11.1,-7.0) -- node[above,font=\scriptsize] {evidence record} (14.8,-7.0);
\node[font=\scriptsize,fill=white,draw=cLine,rounded corners=2pt,text width=3.0cm,align=center] at (14.8,-8.0) {Platt scaling, Trust Score, verdict, sentences};
\draw[darr] (14.8,-9.0) -- node[above,font=\scriptsize] {result} (3.7,-9.0);
\draw[darr] (3.7,-10.0) -- node[above,font=\scriptsize] {card, heatmap, report} (0,-10.0);
\end{tikzpicture}}

  \caption{Sequence diagram of one image check.}\label{fig:seq}
\end{figure}

\section{Class Diagram}
The main software components and their relations are shown in Fig.~\ref{fig:class}. All analysers implement one interface and return an evidence record; the image analyser owns the CLIP encoder and the linear head; the trust engine turns the evidence into a score and a verdict, and the report builder writes the explanation.

\begin{figure}[htbp]
  \centering
  % Class diagram of the main software components (names follow the prototype code)
\resizebox{\linewidth}{!}{%
\begin{tikzpicture}[uml/.style={rectangle split,rectangle split parts=3,draw=cLine,thick,fill=cBlue,font=\footnotesize,align=left,text width=3.6cm,inner sep=3pt}]
\node[uml] (router) at (0,0) {\textbf{ContentRouter}\nodepart{second}type, size limit\nodepart{third}validate(file)\\route(file)};
\node[uml,fill=cGrey] (an) at (6.2,0) {\textbf{Analyser}\;(interface)\nodepart{second}name\nodepart{third}analyse(file): Evidence};
\node[uml,fill=cOrange] (ev) at (12.4,0) {\textbf{Evidence}\nodepart{second}score, visual\\cues, flags\nodepart{third}};
\node[uml,fill=cGreen] (img) at (0,-4.2) {\textbf{ImageAnalyser}\nodepart{second}patches (10), grid\\JPEG quality 95\nodepart{third}analyse(image)};
\node[uml,dashed,fill=white] (vid) at (6.2,-4.2) {\textbf{VideoAnalyser}\;(designed)\nodepart{second}frame rate, face margin\nodepart{third}analyse(video)};
\node[uml,dashed,fill=white] (doc) at (12.4,-4.2) {\textbf{DocumentAnalyser}\;(designed)\nodepart{second}OCR, ELA, SRM, DCT\nodepart{third}analyse(document)};
\node[uml] (clip) at (0,-8.4) {\textbf{ClipEncoder}\nodepart{second}ViT-L/14 (frozen)\nodepart{third}encode(patches)};
\node[uml] (head) at (4.2,-8.4) {\textbf{LinearHead}\nodepart{second}$\mathbf{w}$, $b$, $C$\nodepart{third}logits(features)};
\node[uml,fill=cOrange] (trust) at (8.4,-8.4) {\textbf{TrustEngine}\nodepart{second}$a$, $c$, bands (54, 17)\nodepart{third}trust\_score(p)\\verdict(T)};
\node[uml,fill=cOrange] (rep) at (12.6,-8.4) {\textbf{ReportBuilder}\nodepart{second}evidence, metadata\nodepart{third}explain()\\report()};
\draw[arr] (router) -- node[above,font=\scriptsize] {calls} (an);
\draw[arr] (an) -- node[above,font=\scriptsize] {returns} (ev);
\draw[-{Triangle[open]},thick,cLine] (img.north) -- ++(0,0.9) -| ([xshift=-0.8cm]an.south);
\draw[-{Triangle[open]},thick,cLine] (vid.north) -- (an.south);
\draw[-{Triangle[open]},thick,cLine] (doc.north) -- ++(0,0.9) -| ([xshift=0.8cm]an.south);
\draw[arr] (img.south) -- ++(0,-0.6) -| (clip.north);
\draw[arr] (img.south) -- ++(0,-0.6) -| (head.north);
\draw[arr] (rep.east) -- ++(0.9,0) |- (ev.east);
\draw[arr] (rep.west) -- (trust.east);
\end{tikzpicture}}

  \caption{Class diagram of the main components (dashed classes are designed).}\label{fig:class}
\end{figure}

\section{Screening Procedure}
The screening procedure is shown as a flow chart in Fig.~\ref{fig:flow} and as steps below.

\begin{enumerate}
  \item \textbf{validate}(file): read the type from the bytes, check the size, decode; reject the file if it is corrupt.
  \item \textbf{route}(type): choose the image, video or document analyser.
  \item \textbf{analyse}(file): the analyser returns an evidence record (score, visual, cues, flags).
  \item Calibrate the score into $p$ and compute $T=\mathrm{round}(100(1-p))$.
  \item Assign the verdict from the band of $T$.
  \item Build the explanation from the evidence and show the metadata separately.
\end{enumerate}

\begin{figure}[htbp]
  \centering
  % Flow chart of the screening procedure
\resizebox{!}{0.66\textheight}{%
\begin{tikzpicture}[node distance=0.9cm]
\node[box,rounded corners=9pt,fill=cGrey,text width=3cm] (s) at (0,0) {Start: user uploads a file};
\node[box,text width=3.4cm] (v) at (0,-1.9) {Read file type from the bytes, check size, decode};
\node[diamond,draw=cLine,thick,fill=cOrange,aspect=2.2,align=center,inner sep=1pt,font=\small] (ok) at (0,-4.3) {Valid file?};
\node[box,text width=3.0cm] (err) at (5.2,-4.3) {Show a clear error message};
\node[box,text width=3.4cm] (r) at (0,-6.7) {Route to the image, video or document analyser};
\node[boxg,text width=3.4cm] (a) at (0,-8.6) {Analyser returns an evidence record (score, visual, cues, flags)};
\node[box,text width=3.4cm] (c) at (0,-10.5) {Calibrate the score and compute $T=\mathrm{round}(100(1-p))$};
\node[diamond,draw=cLine,thick,fill=cOrange,aspect=2.4,align=center,inner sep=1pt,font=\small] (b) at (0,-12.9) {Band of $T$};
\node[boxg,text width=2.6cm] (l1) at (-4.6,-15.0) {Likely Authentic};
\node[box,text width=2.6cm] (l2) at (0,-15.0) {Needs Human Review};
\node[boxo,text width=2.6cm] (l3) at (4.6,-15.0) {Likely AI-generated / Manipulated};
\node[box,rounded corners=9pt,fill=cGrey,text width=4.6cm] (e) at (0,-17.0) {Explanation, metadata (shown apart from the score), downloadable report};
\draw[arr] (s) -- (v);
\draw[arr] (v) -- (ok);
\draw[arr] (ok) -- node[right,font=\scriptsize] {yes} (r);
\draw[arr] (ok) -- node[above,font=\scriptsize] {no} (err);
\draw[arr] (r) -- (a);
\draw[arr] (a) -- (c);
\draw[arr] (c) -- (b);
\draw[arr] (b.south west) -- ++(0,-0.4) -| (l1.north);
\draw[arr] (b.south) -- (l2.north);
\draw[arr] (b.south east) -- ++(0,-0.4) -| (l3.north);
\draw[arr] (l1.south) -- ++(0,-0.35) -| (e.north west);
\draw[arr] (l2.south) -- (e.north);
\draw[arr] (l3.south) -- ++(0,-0.35) -| (e.north east);
\end{tikzpicture}}

  \caption{Flow chart of the screening procedure.}\label{fig:flow}
\end{figure}

\section{Additional Checks in the Application}
Three checks sit beside the model and never change the probability it computes.
\begin{itemize}
  \item \textbf{Provenance markers.} Camera data (EXIF) and C2PA credentials \cite{c2pa} are displayed, but they are not scored, because they are easy to remove or forge.
  \item \textbf{Visible Gemini watermark.} Gemini's own image downloads carry a small four-point star near the bottom-right corner. The application matches the corner of the image against a template of this star; when it is found, the alert states that the image was made or edited with Gemini tools, and the verdict is set to AI-generated. The mark can be cropped away, so its absence proves nothing, and it is reported separately from the model's accuracy.
  \item \textbf{Notices for unusual files.} Drawings, plain graphics and screenshots, very small images and files that look like copies re-saved by a messaging app are outside what the detector was built for. The application shows a notice for them; for drawings and graphics it also sets the verdict to \emph{Needs Human Review}.
\end{itemize}

\section{User Interface Design}
The interface is a single web page with drag-and-drop upload. The result card shows the Trust Score in large type with a colour for the verdict (green, amber or red), then the explanation (``Why?''), and below it tabs for the heatmap, the frequency spectrum, error-level analysis, file metadata and technical details. Two further pages show the cross-generator results and describe how the system works. Screenshots are given in Chapter~5.
